\clearpage
\VerdierPhysicalPageAnchor{0201}%
\begingroup\small\setlength{\parskip}{0pt}\setlength{\emergencystretch}{1.5em}
\setlength{\abovedisplayskip}{4pt plus 2pt minus 2pt}%
\setlength{\belowdisplayskip}{4pt plus 2pt minus 2pt}%

\VerdierSemanticAnchor{ch3d-2-3-3-proof-n0-epimorphism}
où $Z^{-q}$ est un élément de $M'$ et les $L^j$,
$-q+1\leq j\leq n_0$, des éléments de $M$. Il existe par suite un
épimorphisme :
\[
\alpha\colon L^{-q}\longrightarrow Z^{-q},
\qquad\text{où }L^{-q}\in M.
\]
Posons alors :
\[
X_{n_0,q}=X_{n_0,q-1}(\alpha,-q),
\qquad
v_{n_0,q}=p(\alpha,-q)\quad\text{(2.1.1)}.
\]

\VerdierSemanticAnchor{ch3d-2-3-3-proof-n0-completion}
Il est clair que par construction on a les propriétés
$I)_{n_0,q}$ et $III)_{n_0,q}$. La propriété
$II)_{n_0,q}$ résulte de ce que l'ensemble $M'$ possède les
propriétés $(E_1)$ et $(E_4)$. On peut donc poursuivre la construction des
$X_{n_0,i}$. En posant
$X_{n_0}=\varprojlim_i X_{n_0,i}$ et en prenant pour
$u_{n_0}\colon X_{n_0}\longrightarrow Y$ le morphisme canonique, on
obtient un complexe et un morphisme qui possèdent les propriétés
$1)_{n_0}$, $2)_{n_0}$, $3)_{n_0}$, $4)_{n_0}$.

\VerdierSemanticAnchor{ch3d-2-3-3-proof-outer-induction}
Pour construire alors les complexes $X_i$ et les morphismes
$u_i\colon X_i\longrightarrow X_{i-1}$, nous procédons par récurrence.
Supposons donc qu'on ait construit pour tout entier $i$, $n_0\leq i<q$,
des complexes $X_i$ et des morphismes
$u_i\colon X_i\longrightarrow X_{i-1}$ possédant les propriétés
$1)_i$, $2)_i$, $3)_i$, $4)_i$, $5)_i$, $n_0\leq i<q$. Il s'agit de
construire un complexe $X_q$ et un morphisme
$u_q\colon X_q\longrightarrow X_{q-1}$ tels que l'on ait les propriétés
$1)_q$, $2)_q$, $3)_q$, $4)_q$, $5)_q$. Tout d'abord, si $Y^q\in M$, on
pose $X_q=X_{q-1}$ et $u_q=\operatorname{id}$. Supposons que $Y^q$ ne
soit pas un élément de $M$, i.e.\ que $M\text{-}\dim(Y^q)>0$.

\VerdierSemanticAnchor{ch3d-2-3-3-proof-q-inner-induction}
Pour construire le complexe $X_q$ et le morphisme
$u_q\colon X_q\longrightarrow X_{q-1}$, nous procédons à nouveau par
récurrence. Soit $r$ un entier tel que
\[
-q<r\leq -q+M\text{-}\dim(Y^q).
\]
Supposons qu'on ait déterminé une suite de complexes $X_{q,l}$,
$-q<l<r$, et de morphismes
$v_{q,l}\colon X_{q,l}\longrightarrow X_{q,l-1}$ (on pose
$X_{q,-q}=X_{q-1}$) tels que :

\VerdierSemanticAnchor{ch3d-2-3-3-proof-q-r-conditions}
\noindent$I')_l$
$(X_{q,l})^j=(X_{q,l-1})^j$, pour tout $j\neq-l,-l+1$.

\noindent$II')_l$
$(X_{q,l})^{-l}\in M'$ et
$M\text{-}\dim((X_{q,l})^{-l})=M\text{-}\dim(Y^q)-l-q$.

\noindent$III')_l$
$(X_{q,l})^{-l+1}\in M$.

\noindent$IV')_l$
$v_{q,l}$ soit un morphisme élémentaire de hauteur $-l+1$ et
$(v_{q,l})^j=\operatorname{id}$, pour $j\neq-l,-l+1$.

\VerdierSemanticAnchor{ch3d-2-3-3-proof-q-rminus1-complex}
Le complexe $X_{q,r-1}$ est alors de la forme :
\[
\cdots\longrightarrow L^{-r}\longrightarrow Z^{-r+1}
\longrightarrow L^{-r+2}\longrightarrow\cdots\longrightarrow L^q
\longrightarrow Y^{q+1}\longrightarrow Y^{q+2}\longrightarrow\cdots,
\]
où les $L^i$ sont des éléments de $M$ et où $Z^{-r+1}$ est un élément de
$M'$ tel que
$M\text{-}\dim(Z^{-r+1})=M\text{-}\dim(Y^q)-r+1-q$. Il existe donc un
épimorphisme
\endgroup
